\documentclass[a4paper]{article}
\usepackage[ urlcolor = blue, colorlinks = true, citecolor = black, linkcolor = black]{hyperref}
\usepackage{graphicx}
%\usepackage{xltxtra}
\usepackage{fancyhdr}
\usepackage{booktabs}
\usepackage{indentfirst}
\usepackage{abstract}

%%%%%%%%%%%%%%%%%%%%%%%%%%%%%%%%%%%%%%%%%%%%%%%%%%%%%%%%%%%%%%%%%%%%%%%%%
% Must use CM font by default, and don't change math font
\usepackage[cm-default,no-math,quiet]{fontspec}
% For using Chinese
\usepackage[SlantFont,CJKnumber]{xeCJK}
\usepackage{CJKnumb}
\punctstyle{banjiao}

% Main Font is Song
\setCJKmainfont[Mapping=tex-text]{song}

% Main English Text is Times New Roman
\setmainfont[Mapping=tex-text]{Times New Roman}



% For same settings as M$ Word.
\usepackage[top=2.5cm,left=2.5cm,bottom=2cm,right=2cm]{geometry}

% For using bibtex
\usepackage{natbib,bibentry}

%do not breaking any words
\usepackage[none]{hyphenat}
\sloppy

% This is for making thebibliography the same style as my school requests
%%%%%%%%%%%%%%%%%%%%%%%%%%%%%%%%%%%%%%%%%%%%%%%%%%%%%%%%%%%%%%%%%%%%%%%%%
\makeatletter\renewenvironment{thebibliography}[1]{%
      \list{\@biblabel{\@arabic\c@enumiv}}%
           {\settowidth\labelwidth{\@biblabel{#1}}%
            \leftmargin\labelwidth%
            \advance\leftmargin\labelsep%
            \@openbib@code%
            \usecounter{enumiv}%
            \let\p@enumiv\@empty%
            \renewcommand\theenumiv{\@arabic\c@enumiv}}%
      \sloppy%
      \clubpenalty4000%
      \@clubpenalty \clubpenalty%
      \widowpenalty4000%
      \sfcode`\.\@m}%
     {\def\@noitemerr%
       {\@latex@warning{Empty `thebibliography' environment}}%
      \endlist}%
\makeatother
%%%%%%%%%%%%%%%%%%%%%%%%%%%%%%%%%%%%%%%%%%%%%%%%%%%%%%%%%%%%%%%%%%%%%%%%%
% new env and cmd
\newlength{\parwidth}%
\newcommand{\framedparbox}[1]{ {% 
 \par 
 \setlength{\parwidth}{\linewidth}%
 \addtolength{\parwidth}{-1.5\parindent}%
% \hspace*{\parindent}% 
 \fbox{%
   \parbox{\parwidth}{\tt{}#1}%
 }%
}\vspace*{1ex}}%

\newenvironment{CJKenumerate}{%
\newcounter{listcount}%
\begin{list}{\stepcounter{listcount}\CJKnumber{\arabic{listcount}}、}{}%
}{\end{list}}

%%%%%%%%%%%%%%%%%%%%
\renewenvironment{abstract}{%
\begin{center}%
{\fontsize{16}{16}\selectfont\bf%三号，加粗，居中
\abstractname}%
\end{center}%
\vspace*{3ex}%
\begin{onehalfspace}%
\par\indent%
}{\end{onehalfspace}}%

\newcommand{\keywordsname}{Keywords:}

\newcommand{\keywords}[1]{\vspace*{3ex}\noindent{\fontsize{14}{14}\selectfont\bf\keywordsname}~#1}
%%%%%%%%%%%%%%%

\newenvironment{centerlastline}{\leftskip0pt plus.5fil \rightskip0pt plus-.5fil \parfillskip0pt plus1fil \relax\parindent 2em\par}{\par}

%%%%%%%%%%%%%%%%%%%%%%%%%%%%%%%%%5

%！！重要！！%这是防止 - -- --- 等特殊符号失去在 LaTeX 下原油的意义。
\defaultfontfeatures{Mapping=tex-text}

% xeCJK 会处理好中文断行
\XeTeXlinebreaklocale{en}
%\XeTeXlinebreakskip=0ex minus 0.5ex plus 1ex 

% 中文日期
\newcommand{\CJKtoday}{\number\year 年 \number\month 月 \number\day 日}

% English title of my thesis
\newcommand{\entitle}{Lexical Analyse in English Machine Translation}
% Chinese title of my thesis
\newcommand{\cntitle}{机器翻译中的英语语义分析}

% Author name
\newcommand{\cnauthor}{蔡万钊}
\newcommand{\enauthor}{Microcai}

% school
\newcommand{\myschool}{浙江理工大学外国语学院}
\newcommand{\myclass}{07英语(2)班}
\newcommand{\mystudentno}{G07710219}
\newcommand{\myteacher}{黄海军}



%\newcommand{\LABTP}{{Lexical Analyse Based Translation Program }}

\title{\entitle}
\author{\enauthor}

\begin{document}

\maketitle

\begin{abstract}
\LABTP extends popular compiler model to support translating human languages.

\end{abstract}

\section{Introduction}

As moving to a new age, people tend to need more communication. Filling the gap
between different cultures requires translating countless books, articles and
words. Human resource is always a limitation. Thousands of pioneers have developed
a serious of Mathematical model for doing translation mechanically. 

Much of the previous model employ too much mathematics, focus not enough to the true meaning
of the words. Translate is not a simple mapping from source language to target language.
In order to do perfect translation, the true meaning of the words must be understood. 

\LABTP is based on the compiler model. What is a compiler? A translator between
one computer language to another. A compiler accept well defined computer language
text, and produce well defined computer language output. The computer language is well
defined, so it is possible to do lexical analyse. As for human languages? Hmm \ldots\ldots 
But, wait! Won't we really do lexical analyse on human language?

Well, It's partly possible. Think of English. English is the mostly
NLP-friendly \footnote{NLP-friendly means, easy to do with Natural Language Processing}
human language ever exist.

\includegraphics{zstu}

\end{document}
